% We formalize a \emph{Compound AI System} as any system comprising one or more coordinated language model invocations, encompassing single-LLM pipelines, multi-agent systems, and Language Programs~\citep{langprobe, better_together, dspy, miprov2}. An LLM-based compound AI system is composed of one or more \textit{language modules}, $M = \langle M_1, \ldots, M_{|M|}\rangle$. A \textit{language module} fully specifies the expected inputs, outputs and behavioural specification of an LLM call made within the compound AI system. A \textit{language module} $M_i$ is fully specified as a tuple $M_i = (\pi_i, \theta_i, IN_i, OUT_i)$, where $\pi_i$ consists of the prompt including instructions and few-shot demonstrations providing a \textit{specification} of the subtask assigned to this module (For example, ``Given \textit{question} as input, generate a \textit{search query} to retrieve documents to help answer this question''), and $\theta_i$ refers to the floating-point weights of the underlying LM that will be used to invoke this subtask. $IN_i$ and $OUT_i$ denote the input and output schemas, respectively, for each call (e.g., $IN_i$ = [question], $OUT_i$ = [search\_query]), and map to concrete values at runtime, for example, $IN_i$ could map to $[question \Rightarrow $ ``What is the capital of France?'' $]$ and $OUT_i$ could map to $[search\_query \Rightarrow $ ``France Capital City'' $]$.

% A Compound AI system is defined as $\Phi = (M, C, IN, OUT)$, with $M = \langle M_1, \ldots, M_{|M|} \rangle$ the set of one or more language modules. The control flow $C$ (specified as code) governs the execution order and conditional logic for invoking modules, orchestrating complex behaviors, interfacing external tools, and applying guardrails. $C$ can perform arbitrary invocations of any of the \emph{language modules}, including invoking them multiple times. $IN$ and $OUT$ specify the global input and output schemas for the system; without loss of generality, $OUT$ always contains a special \texttt{trace} field recording the sequence of module invocations and their respective inputs/outputs.
% A Compound AI System consists of one or more coordinated language model (LLM) invocations, encompassing single-LLM pipelines, multi-agent systems, and language programs~\citep{langprobe, better_together, dspy, miprov2}. Formally, such a system is defined as $\Phi = (M, C, IN, OUT)$, where $M = \langle M_1, \ldots, M_{|M|}\rangle$ is a set of language modules. Each module $M_i = (\pi_i, \theta_i, IN_i, OUT_i)$ specifies $\pi_i$: the prompt (instructions, demonstrations) describing the module’s subtask;
% $\theta_i$: the underlying LLM weights; $IN_i$, $OUT_i$: input/output schemas. The control flow $C$ (typically code) governs the sequencing and conditional logic of module invocation, orchestrating complex behaviors, external tool calls, and guardrails. $C$ may invoke modules arbitrarily, and modules may be called multiple times. The system-level $IN$ and $OUT$ specify the global input/output schemas; $OUT$ always includes a \texttt{trace} of module invocations and respective I/O.

% \paragraph{Compound AI Systems.} 
We follow related work in defining a \textbf{compound AI system} as any modular system composed of one or more language model (LLM) invocations, potentially interleaved with external tool calls, orchestrated through arbitrary control flow. This definition subsumes a broad class of real-world LLM-based AI systems, including \textit{agents}, \textit{multi-agent systems}, and general-purpose scaffolding techniques like ReAct~\citep{react}, Archon~\citep{archon}, etc.
Following~\citet{better_together,dspy,miprov2,langprobe}, we formalize such a system as $\Phi = (M, C, \mathcal{X}, \mathcal{Y})$, where $M = \langle M_1, \ldots, M_{|M|} \rangle$ denotes language modules, $C$ specifies control flow logic, and $\mathcal{X}$, $\mathcal{Y}$ are global input/output schemas. Each module $M_i = (\pi_i, \theta_i, \mathcal{X}_i, \mathcal{Y}_i)$ is an LLM subcomponent: $\pi_i$ is its (system) prompt including instructions and few-shot demonstrations; $\theta_i$ the underlying model weights; $\mathcal{X}_i, \mathcal{Y}_i$ are input/output schemas. At runtime, $C$ orchestrates the sequencing and invocation of modules---e.g., passing outputs from one module to another, invoking modules conditionally, or leveraging tool APIs. This way, $C$ can invoke different modules in any order multiples of times.

% \paragraph{Compound AI System Optimization.}
\label{par:formalizing_cais_optimization}
Given $\Phi$, let $\Pi_\Phi = \langle \pi_1, \ldots, \pi_{|M|}\rangle$ denote the collection of all module prompts and $\Theta_\Phi = \langle \theta_1, \ldots, \theta_{|M|}\rangle$ the set of module weights. The learnable parameters are thus $\langle \Pi, \Theta \rangle_\Phi$. For a task instance $(x, m)$---where $x$ maps to the input schema $\mathcal{X}$ and $m$ contains evaluator metadata (e.g., gold answers, evaluation rubrics, code unit tests)---the system induces an output $y = \Phi(x; \langle \Pi, \Theta \rangle_\Phi )$. A metric $\mu:\mathcal{Y} \times \mathcal{M} \to [0,1]$ then measures the output quality of $y$ with respect to metadata $m$ (for example by calculating, exact match, F1, pass rate, etc.). The optimization problem is thus defined as follows, where $\mathcal{T}$ is a task distribution.:
\begin{align}
    \langle \Pi^*, \Theta^* \rangle_\Phi = \arg\max_{\langle \Pi, \Theta \rangle_\Phi} \mathbb{E}_{(x, m) \sim \mathcal{T}} \left[ \mu\big( \Phi(x; \langle \Pi, \Theta \rangle_\Phi),\, m \big) \right].
\end{align}

\added{We adopt this general problem formulation, allowing updates to both prompts and weights of language modules, to enable comparisons between optimization algorithms that operate in different parameter spaces (e.g., GEPA vs. GRPO).}
% \added{The general problem setting allowing updates to both the prompt and weight of  language modules allow for comparison between optimization algorithms that operate across these parameter spaces.}

\paragraph{Sample-Efficient Optimization.}
In many real-world scenarios, rollouts---concretely, invocations of $\Phi$ plus evaluation by $\mu$---are often computationally, monetarily, or timewise expensive. The optimizer is thus limited to at most $B$ rollouts on a dataset $\mathcal{D}_\text{train} = \{ (x, m)_i \}_{i=1}^N$ with full access to $\mu$. The goal is to identify parameters $\langle \Pi^*, \Theta^* \rangle_\Phi$ that maximize held-out performance, subject to not exceeding the rollout budget $B$:
\begin{align}
    \langle \Pi^*, \Theta^* \rangle_\Phi = \arg\max_{\langle \Pi, \Theta \rangle_\Phi} \mathbb{E}_{(x, m) \sim \mathcal{T}} \left[ \mu\big( \Phi(x; \langle \Pi, \Theta \rangle_\Phi),\, m \big) \right], \quad \text{s.t.} \quad \#\text{rollouts} \leq B.
\end{align}
% This formulation captures the core challenge motivating our work: \emph{How can we extract maximal learning signal from every expensive rollout to enable effective adaptation of complex, modular AI systems in low-data or budget-constrained settings?}
The core challenge, then, is: \emph{How do we extract maximal learning signal from every expensive rollout to enable effective adaptation of complex, modular AI systems in low-data or budget-constrained settings?}

%%%%%%%%%%%%%%%%%%%%%%%%%%%%%%%%%%%%%%%%%%%%%%
% A \textbf{Compound AI System} consists of one or more coordinated language model (LLM) invocations, encompassing single-LLM pipelines, multi-agent systems, and language programs~\citep{langprobe, better_together, dspy, miprov2}. Formally, such a system is defined as $\Phi = (M, C, \mathcal{X}, \mathcal{Y})$, where $M = \langle M_1, \ldots, M_{|M|}\rangle$ is a set of language modules. Each module $M_i = (\pi_i, \theta_i, \mathcal{X}_i, \mathcal{Y}_i)$ fully specifies the subcomponent, where $\pi_i$ is the prompt---including instructions and demonstrations specifying the subtask, for example, ``Given a \textit{question} as input, generate a \textit{search query} to retrieve documents to help answer this question''---and $\theta_i$ denotes the underlying LLM weights. The $\mathcal{X}_i$ and $\mathcal{Y}_i$ specify input and output schemas, e.g., $\mathcal{X}_i = [\text{question}]$, $\mathcal{Y}_i = [\text{search\_query}]$, which map to concrete values at runtime (e.g., $\mathcal{X}_i = [\text{question} \Rightarrow \text{``What is the capital of France?''}]$, $\mathcal{Y}_i = [\text{search\_query} \Rightarrow \text{``France Capital City''}]$). The control flow $C$ (typically implemented in code) governs the sequencing and conditional logic of module invocation, orchestrating complex behaviors, external tool calls, and guardrails---for instance, $C$ might first invoke a question rewriting module, then pass its output to a retrieval module, before finally sending the retrieved information to an answer synthesis module; $C$ may invoke modules multiple times and in arbitrary orders as needed. The system-level $\mathcal{X}$ and $\mathcal{Y}$ specify global input/output schemas, where $\mathcal{Y}$ always includes a \texttt{trace} of module invocations and respective I/O.
% \subsection{Problem Statement: Compound AI System Optimization}
% Given a compound AI system $\Phi = (M, C, \mathcal{X}, \mathcal{Y})$, let $\Pi_\Phi = \langle\pi_1, \ldots, \pi_{|M|}\rangle$ denote the set of prompts, and $\Theta_\Phi = \langle\theta_1, \ldots, \theta_{|M|}\rangle$ denote the set of all weight parameters across modules. Collectively, $\langle \Pi, \Theta \rangle_{\Phi}$ denote the learnable parameters of the compound AI system. A task instance for the system $\Phi$ is specified as a pair $(x, m)$ where $x$ maps to the input schema $\mathcal{X}$ and $m$ contains all the \textit{metadata} necessary to evaluate the task instance (ground truth labels, or unit tests). Let $y = \Phi(x; \langle \Pi, \Theta \rangle_{\Phi})$ denote the output of the system running with parameters $\langle \Pi, \Theta \rangle_{\Phi}$ for input $x$. Then let $\mu : \mathcal{Y} \times \mathcal{M} \to [0,1]$ be an evaluation metric (e.g., exact match, F1, unit tests pass rate, etc.) measuring the quality of the output $y$ produced by the system relative to $m$.

% The compound AI system optimization problem is to find
% \begin{align}
%     \langle \Pi^*, \Theta^* \rangle_\Phi = \arg\max_{\langle \Pi, \Theta \rangle_\Phi} \mathbb{E}_{(x, m) \sim \mathcal{T}} \left[ \mu\left( \Phi(x; \langle \Pi, \Theta \rangle_\Phi),\; m \right) \right],
% \end{align}
% where $\mathcal{T}$ is a distribution over task instances.

% In practise, in domains with expensive rollouts, the optimizer is provided with the system $\Phi$, a training set $\mathcal{D}_\mathrm{train} = \{(x^{(i)}, m^{(i)})\}_{i=1}^N$, access to the evaluation metric $\mu$, and a fixed rollout budget $B$ governing the maximum number of allowed calls to $\mu$. Each query to $\mu$ incurs the cost of a full rollout on a given task instance. The optimizer must output parameters $\langle\Pi^*, \Theta^*\rangle_\Phi$ that approximately maximize expected performance over $\mathcal{T}$, subject to not exceeding the budget $B$ during optimization:
% \[
% \langle \Pi^*, \Theta^* \rangle_\Phi = \arg\max_{\langle\Pi, \Theta\rangle_\Phi} \mathbb{E}_{(x, m) \sim \mathcal{T}}\left[\mu\left( \Phi(x ; \langle\Pi, \Theta\rangle_\Phi),~ m \right)\right].
% \]
% This formulation highlights the central challenge: extract maximal improvement from each costly rollout, enabling efficient optimization of complex AI pipelines when rollouts are expensive.
%%%%%%%%%%%%%%%%%%%%%%%%%%%%%%%%%%%%%%%%%%%%

% consider a dataset $\mathcal{D} = \{(x_i, m_i)\}_{i=1}^N$, where each $x_i$ is an input and $m_i$ contains associated metadata (labels, reference facts, or evaluation resources). 

% The overall objective of compound system optimizer is to maximize the quality of outputs produced by the system according to the evaluation metric $\mu$, in expectation over the input-task distribution. Formally, we seek to optimize the learnable parameters—namely, the collection of prompt specifications $\Pi$ and underlying LLM weights $\Theta$—to maximize performance on the domain of interest. This can be expressed as the following argmax objective:
% \[
% \max_{\Pi,, \Theta} ; \mathbb{E}_{(x, m) \sim \mathcal{D}} , \left[, \mu\left(\Phi(x;, \Pi, \Theta),; m \right) ,\right]
% \]
% where $\Phi(x; \Pi, \Theta)$ denotes the output of the instantiated compound system on input $x$, parameterized by prompts $\Pi$ and weights $\Theta$.

% In practice, we empirically estimate this expectation by evaluating over the dataset $\mathcal{D}$, yielding the optimization problem:
% \[
% \max_{\Pi,, \Theta} ; \frac{1}{|\mathcal{D}|} \sum_{(x, m) \in \mathcal{D}} \mu\left(\Phi(x; \Pi, \Theta),; m\right)
% \]

% This formulation subsumes both prompt optimization (where only $\Pi$ is updated, keeping $\Theta$ fixed) and full model fine-tuning (jointly updating $\Pi$ and $\Theta$), and it readily extends to settings with additional constraints or regularization terms on system behavior.
% We adopt a broad definition of Compound AI System, that encompasses \textit{agents}, \textit{multi-agent systems}, \textit{Language Program}~\citep{langprobe} as well as simple, one-LLM call LLM systems. The following problem specification is in line with previous work including ~\citet{better_together, langprobe, dspy, miprov2}.

% An LLM-based compound AI system is composed of one or more \textit{language modules}, $M = \langle M_1, \ldots, M_{|M|}\rangle$. Each such \textit{language module} fully specifies the expected inputs, outputs and behavioural specification of an LLM call made within the compound AI system. A \textit{language module} $M_i$ is fully specified as a tuple $M_i = (\pi_i, \theta_i, IN_i, OUT_i)$, where $\pi_i$ consists of the prompt including instructions and few-shot demonstrations providing a \textit{specification} of the subtask assigned to this module (For example, ``Given \textit{question} as input, generate a \textit{search query} to retrieve documents to help answer this question''), and $\theta_i$ refers to the floating-point weights of the underlying LM that will be used to invoke this subtask. Finally, $IN_i$ and $OUT_i$ are lists of descriptions of the inputs provided to this LM call (For example, $[question]$), and the expected outputs (For example, $[search\_query]$). While $\pi_i$ and $\theta_i$ pertain to the learnable parameters of the module $M_i$, 

% Now, a compound AI system is specified as a tuple $\Phi = (M, C, IN, OUT)$, where $M$ is the list of one or more \textit{languag modules} as defined above. During the execution of one rollout of the compound AI system, each of the language modules can be invokved one or more times. $C$ refers to the control flow algorithm, specified as arbitrary code, that determines the order of execution of the compound AI system, including all calls to the language modules in $M$. $C$ also incorporates calls to external tools, retrievers, search engines, and developer-defined guardrails, including backtracking. Finally, $IN$ and $OUT$ are lists of input and outputs produced by the compound AI system as a whole (For example, for the multi-hop QA system, $IN = [question]$, and $OUT = [answer]$). $OUT$ additionally always consists of one additional field, $trace$, which is a chronological list of language module invocations, each invocation consisting of $(M_i, in_i, out_i)$, where $M_i$ is the invoked module, $in_i$ and $out_i$ are mappings corresponding to $IN_i$ and $OUT_i$ respectively, 

% \subsection{Problem Statement: Optimizing Compound AI Systems}
% For a compound AI system $\Phi = (M, C, IN, OUT)$, we assume a given dataset $\mathcal{D} = \{(in_1, m_1), \ldots, (in_{N}, m_{N})\}$, where each $in_i$ is a system input (mapping to $IN$) and $m_i$ is \textit{metadata} associated with $in_i$. The metadata $m_i$ provides all the information necessary to evaluate an output $out_i$ produced by the system---for example, a ground truth label, a set of gold-standard supporting facts, or a suite of unit tests for code generation tasks.

% Additionally, we are given a metric $\mu: (OUT, \mathcal{M}) \rightarrow [0, 1]$, which measures the quality or correctness of a system output $out$ given the corresponding metadata $m$ (e.g., exact match, F1, pass@k). For example, for question answering, $m = $ \texttt{[answer = "Paris"]}; $\mu$ compares the generated answer to the label. For information retrieval, $\mu$ can evaluate the fraction of correct documents retrieved, and so on.

% Given such a setup, the goal of the optimization problem is to find the assignment of the prompts and weights of each of the \textit{language modules} in the system, in order to maximize the metric $\mu$ over a problem domain. 

% We formalize a \emph{Compound AI System} as any system comprising one or more coordinated language model invocations, encompassing single-LLM pipelines, multi-agent systems, and Language Programs~\citep{langprobe, better_together, dspy, miprov2}.

% \textbf{Language modules.} The core building block is a \emph{language module} $M_i = (\pi_i, \theta_i, IN_i, OUT_i)$, where $\pi_i$ is a prompt specification (including instructions and demonstrations), and $\theta_i$ are the underlying LLM weights. $IN_i$ and $OUT_i$ denote the input and output schemas, respectively, for each call (e.g., $IN_i$ = [question], $OUT_i$ = [search\_query]).

% \textbf{Compound AI system.} A system is defined as $\Phi = (M, C, IN, OUT)$, with $M = \langle M_1, \ldots, M_{|M|} \rangle$ the set of language modules. The control flow $C$ (specified as code) governs the execution order and conditional logic for invoking modules, orchestrating complex behaviors, interfacing external tools, and applying guardrails. $C$ can perform arbitrary invocations of any of the \textit{language modules}. $IN$ and $OUT$ specify the global input and output schemas for the system; without loss of generality, $OUT$ always contains a special \texttt{trace} field recording the sequence of module invocations and their respective inputs/outputs.

% \subsection{Problem Statement: System Optimization}

% Given a compound system $\Phi = (M, C, IN, OUT)$, consider a dataset $\mathcal{D} = \{(x_i, m_i)\}_{i=1}^N$, where each $x_i$ is an input and $m_i$ contains associated metadata (labels, reference facts, or evaluation resources). Let $\mu : OUT \times \mathcal{M} \to [0,1]$ be an evaluation metric (e.g., exact match, F1, success@k) measuring the quality of a produced output relative to $m$. The objective is to jointly optimize the prompts $\{\pi_i\}$ and model parameters $\{\theta_i\}$ of all modules in order to maximize $\mathbb{E}_{(x, m)\sim\mathcal{D}} \; [\, \mu(\mathrm{Run}_\Phi(x), m)\, ]$ across the task distribution, where $\mathrm{Run}_\Phi(x)$ denotes the system output on input $x$.